% Options for packages loaded elsewhere
\PassOptionsToPackage{unicode}{hyperref}
\PassOptionsToPackage{hyphens}{url}
\documentclass[
]{book}
\usepackage{xcolor}
\usepackage{amsmath,amssymb}
\setcounter{secnumdepth}{5}
\usepackage{iftex}
\ifPDFTeX
  \usepackage[T1]{fontenc}
  \usepackage[utf8]{inputenc}
  \usepackage{textcomp} % provide euro and other symbols
\else % if luatex or xetex
  \usepackage{unicode-math} % this also loads fontspec
  \defaultfontfeatures{Scale=MatchLowercase}
  \defaultfontfeatures[\rmfamily]{Ligatures=TeX,Scale=1}
\fi
\usepackage{lmodern}
\ifPDFTeX\else
  % xetex/luatex font selection
\fi
% Use upquote if available, for straight quotes in verbatim environments
\IfFileExists{upquote.sty}{\usepackage{upquote}}{}
\IfFileExists{microtype.sty}{% use microtype if available
  \usepackage[]{microtype}
  \UseMicrotypeSet[protrusion]{basicmath} % disable protrusion for tt fonts
}{}
\makeatletter
\@ifundefined{KOMAClassName}{% if non-KOMA class
  \IfFileExists{parskip.sty}{%
    \usepackage{parskip}
  }{% else
    \setlength{\parindent}{0pt}
    \setlength{\parskip}{6pt plus 2pt minus 1pt}}
}{% if KOMA class
  \KOMAoptions{parskip=half}}
\makeatother
\usepackage{color}
\usepackage{fancyvrb}
\newcommand{\VerbBar}{|}
\newcommand{\VERB}{\Verb[commandchars=\\\{\}]}
\DefineVerbatimEnvironment{Highlighting}{Verbatim}{commandchars=\\\{\}}
% Add ',fontsize=\small' for more characters per line
\usepackage{framed}
\definecolor{shadecolor}{RGB}{248,248,248}
\newenvironment{Shaded}{\begin{snugshade}}{\end{snugshade}}
\newcommand{\AlertTok}[1]{\textcolor[rgb]{0.94,0.16,0.16}{#1}}
\newcommand{\AnnotationTok}[1]{\textcolor[rgb]{0.56,0.35,0.01}{\textbf{\textit{#1}}}}
\newcommand{\AttributeTok}[1]{\textcolor[rgb]{0.13,0.29,0.53}{#1}}
\newcommand{\BaseNTok}[1]{\textcolor[rgb]{0.00,0.00,0.81}{#1}}
\newcommand{\BuiltInTok}[1]{#1}
\newcommand{\CharTok}[1]{\textcolor[rgb]{0.31,0.60,0.02}{#1}}
\newcommand{\CommentTok}[1]{\textcolor[rgb]{0.56,0.35,0.01}{\textit{#1}}}
\newcommand{\CommentVarTok}[1]{\textcolor[rgb]{0.56,0.35,0.01}{\textbf{\textit{#1}}}}
\newcommand{\ConstantTok}[1]{\textcolor[rgb]{0.56,0.35,0.01}{#1}}
\newcommand{\ControlFlowTok}[1]{\textcolor[rgb]{0.13,0.29,0.53}{\textbf{#1}}}
\newcommand{\DataTypeTok}[1]{\textcolor[rgb]{0.13,0.29,0.53}{#1}}
\newcommand{\DecValTok}[1]{\textcolor[rgb]{0.00,0.00,0.81}{#1}}
\newcommand{\DocumentationTok}[1]{\textcolor[rgb]{0.56,0.35,0.01}{\textbf{\textit{#1}}}}
\newcommand{\ErrorTok}[1]{\textcolor[rgb]{0.64,0.00,0.00}{\textbf{#1}}}
\newcommand{\ExtensionTok}[1]{#1}
\newcommand{\FloatTok}[1]{\textcolor[rgb]{0.00,0.00,0.81}{#1}}
\newcommand{\FunctionTok}[1]{\textcolor[rgb]{0.13,0.29,0.53}{\textbf{#1}}}
\newcommand{\ImportTok}[1]{#1}
\newcommand{\InformationTok}[1]{\textcolor[rgb]{0.56,0.35,0.01}{\textbf{\textit{#1}}}}
\newcommand{\KeywordTok}[1]{\textcolor[rgb]{0.13,0.29,0.53}{\textbf{#1}}}
\newcommand{\NormalTok}[1]{#1}
\newcommand{\OperatorTok}[1]{\textcolor[rgb]{0.81,0.36,0.00}{\textbf{#1}}}
\newcommand{\OtherTok}[1]{\textcolor[rgb]{0.56,0.35,0.01}{#1}}
\newcommand{\PreprocessorTok}[1]{\textcolor[rgb]{0.56,0.35,0.01}{\textit{#1}}}
\newcommand{\RegionMarkerTok}[1]{#1}
\newcommand{\SpecialCharTok}[1]{\textcolor[rgb]{0.81,0.36,0.00}{\textbf{#1}}}
\newcommand{\SpecialStringTok}[1]{\textcolor[rgb]{0.31,0.60,0.02}{#1}}
\newcommand{\StringTok}[1]{\textcolor[rgb]{0.31,0.60,0.02}{#1}}
\newcommand{\VariableTok}[1]{\textcolor[rgb]{0.00,0.00,0.00}{#1}}
\newcommand{\VerbatimStringTok}[1]{\textcolor[rgb]{0.31,0.60,0.02}{#1}}
\newcommand{\WarningTok}[1]{\textcolor[rgb]{0.56,0.35,0.01}{\textbf{\textit{#1}}}}
\usepackage{longtable,booktabs,array}
\newcounter{none} % for unnumbered tables
\usepackage{calc} % for calculating minipage widths
% Correct order of tables after \paragraph or \subparagraph
\usepackage{etoolbox}
\makeatletter
\patchcmd\longtable{\par}{\if@noskipsec\mbox{}\fi\par}{}{}
\makeatother
% Allow footnotes in longtable head/foot
\IfFileExists{footnotehyper.sty}{\usepackage{footnotehyper}}{\usepackage{footnote}}
\makesavenoteenv{longtable}
\usepackage{graphicx}
\makeatletter
\newsavebox\pandoc@box
\newcommand*\pandocbounded[1]{% scales image to fit in text height/width
  \sbox\pandoc@box{#1}%
  \Gscale@div\@tempa{\textheight}{\dimexpr\ht\pandoc@box+\dp\pandoc@box\relax}%
  \Gscale@div\@tempb{\linewidth}{\wd\pandoc@box}%
  \ifdim\@tempb\p@<\@tempa\p@\let\@tempa\@tempb\fi% select the smaller of both
  \ifdim\@tempa\p@<\p@\scalebox{\@tempa}{\usebox\pandoc@box}%
  \else\usebox{\pandoc@box}%
  \fi%
}
% Set default figure placement to htbp
\def\fps@figure{htbp}
\makeatother
\setlength{\emergencystretch}{3em} % prevent overfull lines
\providecommand{\tightlist}{%
  \setlength{\itemsep}{0pt}\setlength{\parskip}{0pt}}
\usepackage[]{natbib}
\bibliographystyle{apalike}
% LaTeX preamble for PDF output
\usepackage{booktabs}
\usepackage{longtable}
\usepackage{array}
\usepackage{multirow}
\usepackage{wrapfig}
\usepackage{float}
\usepackage{colortbl}
\usepackage{pdflscape}
\usepackage{tabu}
\usepackage{threeparttable}
\usepackage{threeparttablex}
\usepackage[normalem]{ulem}
\usepackage{makecell}
\usepackage{xcolor}

% ForestPlots green
\definecolor{fpgreen}{HTML}{2d6a4f}

% Hyperlink colour
\hypersetup{colorlinks=true, linkcolor=fpgreen, urlcolor=fpgreen, citecolor=fpgreen}
\usepackage{bookmark}
\IfFileExists{xurl.sty}{\usepackage{xurl}}{} % add URL line breaks if available
\urlstyle{same}
\hypersetup{
  pdftitle={ForestPlots.net Training Manual},
  pdfauthor={ForestPlots Team, University of Leeds},
  hidelinks,
  pdfcreator={LaTeX via pandoc}}

\title{ForestPlots.net Training Manual}
\author{ForestPlots Team, University of Leeds}
\date{2026-10-16}

\begin{document}
\maketitle

{
\setcounter{tocdepth}{1}
\tableofcontents
}
\chapter*{Welcome}\label{welcome}
\addcontentsline{toc}{chapter}{Welcome}

Welcome to the \textbf{ForestPlots.net Training Manual}. This guide is designed to help new and existing users of \href{https://forestplots.net}{ForestPlots.net} --- the global network for monitoring tropical forest plots --- to get the most out of the platform.

\section*{How to use this manual}\label{how-to-use-this-manual}
\addcontentsline{toc}{section}{How to use this manual}

This manual is organised into chapters that follow a natural workflow, from creating an account through to downloading and analysing data. You can read it sequentially or jump to the chapter most relevant to your needs.

\begin{itemize}
\tightlist
\item
  \textbf{New users} should start with Chapters \ref{introduction} and \ref{getting-started}.
\item
  \textbf{Data contributors} will find Chapters \ref{data-entry} and \ref{data-quality} most useful.
\item
  \textbf{Data users} seeking to query or download data should head to Chapter \ref{querying-data}.
\item
  \textbf{R users} will find the \texttt{RAINFOR} / ForestPlots R package covered in Chapter \ref{r-package}.
\end{itemize}

\section*{About ForestPlots.net}\label{about-forestplots.net}
\addcontentsline{toc}{section}{About ForestPlots.net}

ForestPlots.net is a pan-tropical network and database hosting long-term forest inventory data from permanent sample plots across the tropics. It supports research into forest dynamics, carbon stocks, and ecosystem responses to climate change.

\chapter{Introduction}\label{introduction}

\section{What is ForestPlots.net?}\label{what-is-forestplots.net}

\emph{{[}Content to be added{]}}

\section{The science behind the network}\label{the-science-behind-the-network}

\emph{{[}Content to be added{]}}

\section{Key terminology}\label{key-terminology}

{\def\LTcaptype{none} % do not increment counter
\begin{longtable}[]{@{}ll@{}}
\toprule\noalign{}
Term & Definition \\
\midrule\noalign{}
\endhead
\bottomrule\noalign{}
\endlastfoot
Plot & \emph{{[}Definition{]}} \\
Census & \emph{{[}Definition{]}} \\
Stem & \emph{{[}Definition{]}} \\
DBH & \emph{{[}Definition{]}} \\
\end{longtable}
}

\section{Data ownership and governance}\label{data-ownership-and-governance}

\emph{{[}Content to be added{]}}

\section{How to cite ForestPlots.net}\label{how-to-cite-forestplots.net}

If you use data or the platform in your research, please cite:

\begin{quote}
\emph{{[}Insert preferred citation here{]}}
\end{quote}

For data downloaded from specific plots or networks, additional citations may be required. See the Data Download chapter for details.

\chapter{Getting Started}\label{getting-started}

\section{Creating an account}\label{creating-an-account}

\emph{{[}Content to be added{]}}

\section{Logging in}\label{logging-in}

\emph{{[}Content to be added{]}}

\section{Account types and permissions}\label{account-types-and-permissions}

{\def\LTcaptype{none} % do not increment counter
\begin{longtable}[]{@{}ll@{}}
\toprule\noalign{}
Role & Permissions \\
\midrule\noalign{}
\endhead
\bottomrule\noalign{}
\endlastfoot
Viewer & \emph{{[}Description{]}} \\
Contributor & \emph{{[}Description{]}} \\
Plot Owner & \emph{{[}Description{]}} \\
Administrator & \emph{{[}Description{]}} \\
\end{longtable}
}

\section{Joining a network or group}\label{joining-a-network-or-group}

\emph{{[}Content to be added{]}}

\section{Requesting access to specific plots}\label{requesting-access-to-specific-plots}

\emph{{[}Content to be added{]}}

\section{Your user profile}\label{your-user-profile}

\emph{{[}Content to be added{]}}

\chapter{Navigating the Platform}\label{navigating}

\section{The dashboard}\label{the-dashboard}

\emph{{[}Content to be added{]}}

\section{Finding plots}\label{finding-plots}

\emph{{[}Content to be added{]}}

\section{The map interface}\label{the-map-interface}

\emph{{[}Content to be added{]}}

\section{Plot summary pages}\label{plot-summary-pages}

\emph{{[}Content to be added{]}}

\section{Navigation tips and shortcuts}\label{navigation-tips-and-shortcuts}

\emph{{[}Content to be added{]}}

\chapter{Data Entry}\label{data-entry}

\section{Overview of the data entry workflow}\label{overview-of-the-data-entry-workflow}

\emph{{[}Content to be added{]}}

\section{Creating a new plot}\label{creating-a-new-plot}

\emph{{[}Content to be added{]}}

\section{Adding plot metadata}\label{adding-plot-metadata}

\emph{{[}Content to be added{]}}

\section{Creating a census}\label{creating-a-census}

\emph{{[}Content to be added{]}}

\section{Entering stem-level data}\label{entering-stem-level-data}

\subsection{Species identification}\label{species-identification}

\emph{{[}Content to be added{]}}

\subsection{Measurements}\label{measurements}

\emph{{[}Content to be added{]}}

\subsection{Stem status codes}\label{stem-status-codes}

{\def\LTcaptype{none} % do not increment counter
\begin{longtable}[]{@{}ll@{}}
\toprule\noalign{}
Code & Meaning \\
\midrule\noalign{}
\endhead
\bottomrule\noalign{}
\endlastfoot
A & Alive \\
D & Dead \\
R & Recruit \\
M & Missing \\
\end{longtable}
}

\emph{{[}Add further codes as appropriate{]}}

\section{Uploading data in bulk}\label{uploading-data-in-bulk}

\emph{{[}Content to be added{]}}

\section{Editing existing records}\label{editing-existing-records}

\emph{{[}Content to be added{]}}

\chapter{Data Quality and Validation}\label{data-quality}

\section{Why data quality matters}\label{why-data-quality-matters}

\emph{{[}Content to be added{]}}

\section{Automated validation checks}\label{automated-validation-checks}

\emph{{[}Content to be added{]}}

\section{Reviewing validation flags}\label{reviewing-validation-flags}

\emph{{[}Content to be added{]}}

\section{Manual quality review}\label{manual-quality-review}

\emph{{[}Content to be added{]}}

\section{Common data entry errors and how to fix them}\label{common-data-entry-errors-and-how-to-fix-them}

{\def\LTcaptype{none} % do not increment counter
\begin{longtable}[]{@{}lll@{}}
\toprule\noalign{}
Error & Likely cause & How to fix \\
\midrule\noalign{}
\endhead
\bottomrule\noalign{}
\endlastfoot
\emph{{[}Error type{]}} & \emph{{[}Cause{]}} & \emph{{[}Fix{]}} \\
\end{longtable}
}

\section{Data approval and sign-off}\label{data-approval-and-sign-off}

\emph{{[}Content to be added{]}}

\chapter{Querying and Downloading Data}\label{querying-data}

\section{Overview of the query system}\label{overview-of-the-query-system}

\emph{{[}Content to be added{]}}

\section{Searching for data}\label{searching-for-data}

\emph{{[}Content to be added{]}}

\section{Filtering by plot characteristics}\label{filtering-by-plot-characteristics}

\emph{{[}Content to be added{]}}

\section{Downloading data}\label{downloading-data}

\emph{{[}Content to be added{]}}

\section{Understanding the download files}\label{understanding-the-download-files}

\subsection{File structure}\label{file-structure}

\emph{{[}Content to be added{]}}

\subsection{Key columns}\label{key-columns}

{\def\LTcaptype{none} % do not increment counter
\begin{longtable}[]{@{}lll@{}}
\toprule\noalign{}
Column & Description & Units \\
\midrule\noalign{}
\endhead
\bottomrule\noalign{}
\endlastfoot
\texttt{PlotCode} & Unique plot identifier & --- \\
\texttt{CensusNo} & Census number within a plot & --- \\
\texttt{Tag} & Individual stem tag number & --- \\
\texttt{DBH} & Diameter at breast height & cm \\
\texttt{Status} & Stem status at census & --- \\
\end{longtable}
}

\emph{{[}Expand table as needed{]}}

\section{Data use agreements and attribution}\label{data-use-agreements-and-attribution}

\emph{{[}Content to be added{]}}

\chapter{The ForestPlots R Package}\label{r-package}

\section{Installation}\label{installation}

\begin{Shaded}
\begin{Highlighting}[]
\CommentTok{\# Install from CRAN (if available)}
\FunctionTok{install.packages}\NormalTok{(}\StringTok{"forestplots"}\NormalTok{)}

\CommentTok{\# Or install the development version from GitHub}
\CommentTok{\# install.packages("devtools")}
\NormalTok{devtools}\SpecialCharTok{::}\FunctionTok{install\_github}\NormalTok{(}\StringTok{"ForestPlots/forestplots{-}r"}\NormalTok{)}
\end{Highlighting}
\end{Shaded}

\section{Authentication}\label{authentication}

\begin{Shaded}
\begin{Highlighting}[]
\FunctionTok{library}\NormalTok{(forestplots)}

\CommentTok{\# Authenticate with your credentials}
\FunctionTok{fp\_auth}\NormalTok{(}\AttributeTok{username =} \StringTok{"your\_username"}\NormalTok{, }\AttributeTok{password =} \StringTok{"your\_password"}\NormalTok{)}

\CommentTok{\# Or use an API token (recommended)}
\FunctionTok{fp\_auth}\NormalTok{(}\AttributeTok{token =} \StringTok{"your\_api\_token"}\NormalTok{)}
\end{Highlighting}
\end{Shaded}

\section{Loading data into R}\label{loading-data-into-r}

\begin{Shaded}
\begin{Highlighting}[]
\CommentTok{\# Download stem data for a specific plot}
\NormalTok{stems }\OtherTok{\textless{}{-}} \FunctionTok{fp\_get\_stems}\NormalTok{(}\AttributeTok{plot\_code =} \StringTok{"BOL{-}01"}\NormalTok{)}

\CommentTok{\# Download all plots you have access to}
\NormalTok{my\_plots }\OtherTok{\textless{}{-}} \FunctionTok{fp\_get\_plots}\NormalTok{()}
\end{Highlighting}
\end{Shaded}

\section{Common analyses}\label{common-analyses}

\subsection{Calculating basal area}\label{calculating-basal-area}

\begin{Shaded}
\begin{Highlighting}[]
\CommentTok{\# }\AlertTok{TODO}\CommentTok{: Add working example}
\end{Highlighting}
\end{Shaded}

\subsection{Estimating above-ground biomass (AGB)}\label{estimating-above-ground-biomass-agb}

\begin{Shaded}
\begin{Highlighting}[]
\CommentTok{\# }\AlertTok{TODO}\CommentTok{: Add working example using BIOMASS or equivalent}
\end{Highlighting}
\end{Shaded}

\subsection{Computing stem growth and mortality rates}\label{computing-stem-growth-and-mortality-rates}

\begin{Shaded}
\begin{Highlighting}[]
\CommentTok{\# }\AlertTok{TODO}\CommentTok{: Add working example}
\end{Highlighting}
\end{Shaded}

\section{Working with the downloaded CSV files}\label{working-with-the-downloaded-csv-files}

\begin{Shaded}
\begin{Highlighting}[]
\CommentTok{\# Load stem data}
\NormalTok{stems }\OtherTok{\textless{}{-}} \FunctionTok{read.csv}\NormalTok{(}\StringTok{"stems.csv"}\NormalTok{)}

\CommentTok{\# Load plot metadata}
\NormalTok{plots }\OtherTok{\textless{}{-}} \FunctionTok{read.csv}\NormalTok{(}\StringTok{"plots.csv"}\NormalTok{)}

\CommentTok{\# Join tables}
\FunctionTok{library}\NormalTok{(dplyr)}
\NormalTok{data }\OtherTok{\textless{}{-}} \FunctionTok{left\_join}\NormalTok{(stems, plots, }\AttributeTok{by =} \StringTok{"PlotCode"}\NormalTok{)}
\end{Highlighting}
\end{Shaded}

\section{Further resources}\label{further-resources}

\begin{itemize}
\tightlist
\item
  Package documentation: \emph{{[}link{]}}
\item
  GitHub repository: \emph{{[}link{]}}
\item
  Example workflows: \emph{{[}link{]}}
\end{itemize}

\chapter{Advanced Features}\label{advanced-features}

\section{Managing a network or group}\label{managing-a-network-or-group}

\emph{{[}Content to be added{]}}

\section{Plot documents and attachments}\label{plot-documents-and-attachments}

\emph{{[}Content to be added{]}}

\section{Species lists and taxonomic management}\label{species-lists-and-taxonomic-management}

\emph{{[}Content to be added{]}}

\section{API access}\label{api-access}

\subsection{Authentication}\label{authentication-1}

\begin{Shaded}
\begin{Highlighting}[]
\CommentTok{\# Example API token request}
\ExtensionTok{curl} \AttributeTok{{-}X}\NormalTok{ POST https://forestplots.net/api/v1/auth }\DataTypeTok{\textbackslash{}}
  \AttributeTok{{-}H} \StringTok{"Content{-}Type: application/json"} \DataTypeTok{\textbackslash{}}
  \AttributeTok{{-}d} \StringTok{\textquotesingle{}\{"username": "your\_username", "password": "your\_password"\}\textquotesingle{}}
\end{Highlighting}
\end{Shaded}

\subsection{Example API calls}\label{example-api-calls}

\begin{Shaded}
\begin{Highlighting}[]
\CommentTok{\# }\AlertTok{TODO}\CommentTok{: Add R examples using httr or the ForestPlots R package\textquotesingle{}s API layer}
\end{Highlighting}
\end{Shaded}

\section{Exporting data for other platforms}\label{exporting-data-for-other-platforms}

\emph{{[}Content to be added{]}}

\section{Notifications and alerts}\label{notifications-and-alerts}

\emph{{[}Content to be added{]}}

\chapter{Frequently Asked Questions}\label{faq}

\section{Account and access}\label{account-and-access}

\textbf{Q: I haven't received my account verification email --- what should I do?}

\emph{{[}Answer{]}}

\textbf{Q: How do I reset my password?}

\emph{{[}Answer{]}}

\textbf{Q: I can see a plot on the map but cannot access its data --- why?}

\emph{{[}Answer{]}}

\begin{center}\rule{0.5\linewidth}{0.5pt}\end{center}

\section{Data entry}\label{data-entry-1}

\textbf{Q: What should I do if a stem has died between censuses?}

\emph{{[}Answer{]}}

\textbf{Q: Can I enter data from multiple censuses at the same time?}

\emph{{[}Answer{]}}

\textbf{Q: How do I handle a re-tagged stem where the original tag was lost?}

\emph{{[}Answer{]}}

\begin{center}\rule{0.5\linewidth}{0.5pt}\end{center}

\section{Data quality}\label{data-quality-1}

\textbf{Q: I received a validation error for a DBH value I am confident is correct --- what do I do?}

\emph{{[}Answer{]}}

\textbf{Q: Can I override a validation flag?}

\emph{{[}Answer{]}}

\begin{center}\rule{0.5\linewidth}{0.5pt}\end{center}

\section{Querying and downloading}\label{querying-and-downloading}

\textbf{Q: I can see a plot but the download button is greyed out --- why?}

\emph{{[}Answer{]}}

\textbf{Q: How do I request access to data that is under embargo?}

\emph{{[}Answer{]}}

\textbf{Q: What format should I cite when using downloaded data in a publication?}

\emph{{[}Answer{]}}

\begin{center}\rule{0.5\linewidth}{0.5pt}\end{center}

\section{R package}\label{r-package-1}

\textbf{Q: The R package says my credentials are invalid, but I can log in on the website --- what's wrong?}

\emph{{[}Answer{]}}

\textbf{Q: How do I report a bug in the R package?}

\emph{{[}Answer{]}}

\begin{center}\rule{0.5\linewidth}{0.5pt}\end{center}

\section{Getting further help}\label{getting-further-help}

If your question is not answered here, please contact the ForestPlots.net team:

\begin{itemize}
\tightlist
\item
  \textbf{Email:} \emph{{[}support email{]}}
\item
  \textbf{Help forum:} \emph{{[}link{]}}
\item
  \textbf{GitHub Issues (R package):} \emph{{[}link{]}}
\end{itemize}

\bibliography{references.bib,packages.bib}

\end{document}
